\section{Circuit, code and post-selection conventions}\label{app:circuits}
This appendix specifies the circuits, code conventions and the post-selection criteria needed to reproduce the numbers in the main text.
The $d=5$ circuit is the corrected cultivation circuit distributed with \href{https://github.com/unitaryfoundation/clifft/blob/main/tests/fixtures/cultivation_d5.stim}{Clifft} and \href{https://github.com/haoliri0/SymFT_Test/blob/main/benchmark/circuit/msc_d5_inject_cultivate_p1e-3.stim}{SymFT}. The standalone $d=3$ circuit is the \href{https://github.com/unitaryfoundation/clifft/blob/main/docs/guide/circuits/circuit_d3_t_gate_p0.001.stim}{Clifft example}, with its \href{https://github.com/haoliri0/SymFT_Test/blob/main/benchmark/circuit/msc_d3_inject_cultivate_p1e-3.stim}{SymFT counterpart}.\footnote{The exact $d=5$ file used has SHA-256 checksum \nolinkurl{859d925c74712efaae8bab1d944f9e5be4e042e226eeafcc6d7bef9d6fbfbae5}. A copy with this checksum is exactly the circuit analysed, byte-for-byte.} The distributed file carries a fixed noise value in its annotations. For each result we replace it, at every location, by the value of $p$ quoted with that result, taken as an exact fraction (for example $p=1/1000$), so that the fixed-$p$ probabilities are rational numbers.


There are 3,564 fault locations in the full $d=5$ cultivation circuit of Refs.~\cite{clifft,soft}. Exactly 1,568 of them have identity response, meaning that none of their outcomes changes $A$ or $B$, so summing over their outcomes gives $(1-p)+p=1$ and they drop out exactly without being set noiseless. The remaining 1,996 consist of
\begin{itemize}
\item injection: 244,
\item $d=3$ double checking: 158,
\item code growth: 908,
\item $d=5$ double checking: 686.
\end{itemize}
The standalone $d=3$ circuit has 518 fault locations, of which 402 have non-identity response.


\subsection{The colour code and the \texorpdfstring{$T$}{T} layers}
For a colour code with face operators $X_f,Z_f$, the terminal code projector and
wrong-outcome projector are
\begin{equation}
 \begin{gathered}
 \Pi_0=\prod_f\frac{I+X_f}{2}\frac{I+Z_f}{2},\\
 \Pi_{\rm bad}=\Pi_0\frac{I-M_L}{2},\qquad
 M_L=\frac{X_L-Y_L}{\sqrt2}.
 \end{gathered}
 \label{eq:projectors}
\end{equation}
In this circuit the target state is $\ket{T_L^\dagger}=(\ket{0_L}+e^{-i\pi/4}\ket{1_L})/\sqrt2$, the $+1$ eigenstate of $M_L$, and the terminal readout that tests it is noiseless.

For $d=3$, order the data sites as $(0,3,7,9,10,12,13)$ and write the parity-check matrix of the $[[7,1,3]]$ colour code as\footnote{The $X$ and $Z$ parity-check matrices are the same.}:
\begin{equation}
 H^{(3)}=\begin{pmatrix}
 1&1&1&0&1&0&0\\
 0&1&1&1&0&1&0\\
 0&0&1&0&1&1&1
 \end{pmatrix}.
 \label{eq:H3}
\end{equation}
Each row gives the support of both an $X$ and a $Z$ face stabiliser.
If $C$ is the row span of $H$, the logical basis is
\begin{equation}
 \ket{0_L}=|C|^{-1/2}\sum_{c\in C}\ket c,\qquad
 \ket{1_L}=X_{\rm all}\ket{0_L}.
 \label{eq:codebasis}
\end{equation}
Here $X_{\rm all}=\prod_{q}X_q$, the product over all data sites, is the transversal $X$, which is $X_L$ for these codes.
Write an $X$ operator on the data qubits as a mask, the binary vector $\boldsymbol a$ of Eq.~\eqref{eq:address} with $a_q=1$ where the operator acts. In the site order $(0,3,7,9,10,12,13)$, for example, $X_7$ has mask $\boldsymbol a=(0,0,1,0,0,0,0)$, the blue face $X_0X_3X_7X_{10}$ has mask $(1,1,1,0,1,0,0)$, the first row of $H^{(3)}$, and $X_L$ has mask $\boldsymbol1=(1,1,1,1,1,1,1)$. The masks of the $X$ stabilisers form the row span $C$ of $H$, and the masks that no $Z$ face detects form its orthogonal complement $C^\perp$. Every face has even weight and even overlap with every other face, so $C$ contains only even-weight masks, while the all-ones mask $\boldsymbol1$, the support of $X_L$, is orthogonal to every face and has odd weight. Hence $C^\perp=C\oplus\spann\{\boldsymbol1\}$: an undetected $X$ mask is either a stabiliser or a stabiliser times $X_L$, and its weight parity, the bit $\boldsymbol1\cdot\boldsymbol a$ of Eq.~\eqref{eq:address}, tells which. The same holds for $Z$ masks with $Z_L$, since the code is self-dual, and Appendix~\ref{app:response} uses this decomposition again.

\begin{table}[htbp]\centering
\captionsetup{position=above,skip=4pt}
\caption{Face supports defining the $d=5$ incidence matrix.}
\label{tab:faces}
\begin{tabular}{rll}\toprule
Face & Colour & $d=5$ source sites\\\midrule
1 & Blue & $0,3,7,11$\\
2 & Blue & $9,13,15,19,21,27$\\
3 & Blue & $17,25,32,36$\\
4 & Red & $3,7,9,13$\\
5 & Red & $15,21,23,29$\\
6 & Red & $19,25,27,32,34,38$\\
7 & Green & $7,11,13,17,19,25$\\
8 & Green & $21,27,29,34$\\
9 & Green & $32,36,38,40$\\\bottomrule
\end{tabular}
\end{table}

Let $D$ be the exit layer of signed $T$ gates, highlighted in Fig.~\ref{fig:d3circuit}, and $D^\dagger$ the entry layer.
For $d=3$, $D=T^{\otimes7}$. For $d=5$, $D^\dagger$ has $T$ on
sites $0,11,17,19,36,40$ and $T^\dagger$ on the other thirteen data sites,
as drawn in Fig.~\ref{fig:patches}. The encoded check observable is
\begin{equation}
 U=D X_{\rm all}D^\dagger,\qquad K_+=\frac{I+U}{2}.
 \label{eq:idealU}
\end{equation}
Figure~\ref{fig:ideal} is the idealised version of this projector on the code space, with one ancillary qubit and one readout. The double-checking sub-circuits actually used are the CNOT circuits in Figs.~\ref{fig:d3circuit} and~\ref{fig:d5circuit}. The noisy calculation uses those circuits in full, with all their measurement outcomes and fault locations, one fault location at every operation as in the standard circuit-level noise model.

\diagram{ideal_check}{The ideal encoded projector associated with double checking.}{fig:ideal}

\subsection{Double checking circuits}
Figures~\ref{fig:d3circuit} and \ref{fig:d5circuit} show the gate sequences of the cultivation double-checking sub-circuits.
In the $d=3$ double-checking circuit, there are four layers of CNOTs, followed by the
central $X$ measurement$+$reset and then reverse unfolding CNOT circuit and final measurements.
Similarly, the $d=5$ double-checking circuit has seven fold layers with varying number of CNOT(s) in each layer followed by similar fold, unfold and measurements.




These double-checking circuits produce many measured bits and post-selection is a set of modulo 2 addition conditions on these measured bits. 
In the $d=3$ sub-circuit, the central readout bit $r_8$ and the auxiliary readout bits enter conditions such as:
\[r_6\oplus r_7\oplus r_8=0 \ , \]
\[r_{11}\oplus r_8=0 \ \text{etc} \ ,\]
these conditions/constraints are known as \textit{detectors}.
Post-selection acceptance occurs when all these detectors return bit 0. 
However, a fault pattern (consisting of many faults) can flips a logical operator but bypasses these detector constraints, as these faults can collectively cancel out in the $\oplus$ addition. 
This leads to a malignant fault pattern that is not caught by post-selection.
This is how the two-fault logical error of Fig.~\ref{fig:d3-fixed-faults} passes post-selection.

\begin{figure}[htbp]\centering
\begin{adjustbox}{max width=.98\linewidth,max totalheight=.82\textheight}
\begin{tikzpicture}[font=\scriptsize,line width=.45pt]
\begin{yquant}[register/minimum height=1.1mm,register/minimum depth=1.1mm,register/separation=.6mm,operator/minimum width=4mm,operator/separation=2.2mm]
qubit {$6\;\ket+$} q6;
qubit {$7$} q7;
qubit {$8\;\ket+$} q8;
qubit {$9$} q9;
qubit {$2\;\ket+$} q2;
qubit {$3$} q3;
qubit {$1\;\ket+$} q1;
qubit {$0$} q0;
qubit {$11\;\ket+$} q11;
qubit {$10$} q10;
qubit {$14\;\ket+$} q14;
qubit {$13$} q13;
qubit {$12$} q12;
[name=dentry, fill=highlight!8] box {$T^\dagger$} q7,q9,q3,q0,q10,q13,q12;
align q6,q7,q8,q9,q2,q3,q1,q0,q11,q10,q14,q13,q12;
cnot q0 | q1;
cnot q3 | q2;
cnot q7 | q6;
cnot q9 | q8;
cnot q10 | q11;
cnot q13 | q14;
align q6,q7,q8,q9,q2,q3,q1,q0,q11,q10,q14,q13,q12;
cnot q1 | q3;
cnot q8 | q7;
cnot q14 | q11;
align q6,q7,q8,q9,q2,q3,q1,q0,q11,q10,q14,q13,q12;
cnot q3 | q7;
cnot q12 | q11;
align q6,q7,q8,q9,q2,q3,q1,q0,q11,q10,q14,q13,q12;
cnot q11 | q7;
align q6,q7,q8,q9,q2,q3,q1,q0,q11,q10,q14,q13,q12;
[name=m8, type=qubit] measure {$X$} q7;
box {$R_X$} q7;
align q6,q7,q8,q9,q2,q3,q1,q0,q11,q10,q14,q13,q12;
cnot q11 | q7;
align q6,q7,q8,q9,q2,q3,q1,q0,q11,q10,q14,q13,q12;
cnot q3 | q7;
cnot q12 | q11;
align q6,q7,q8,q9,q2,q3,q1,q0,q11,q10,q14,q13,q12;
cnot q1 | q3;
cnot q8 | q7;
cnot q14 | q11;
align q6,q7,q8,q9,q2,q3,q1,q0,q11,q10,q14,q13,q12;
cnot q0 | q1;
cnot q3 | q2;
cnot q7 | q6;
cnot q9 | q8;
cnot q10 | q11;
cnot q13 | q14;
align q6,q7,q8,q9,q2,q3,q1,q0,q11,q10,q14,q13,q12;
[name=dexit, fill=highlight!15] box {$T$} q7,q9,q3,q0,q10,q13,q12;
align q6,q7,q8,q9,q2,q3,q1,q0,q11,q10,q14,q13,q12;
measure {$X$} q14,q11,q6,q2,q8,q1;
output {\textcolor{highlight}{$r_{9}$}} q14;
output {\textcolor{highlight}{$r_{10}$}} q11;
output {\textcolor{highlight}{$r_{11}$}} q6;
output {\textcolor{highlight}{$r_{12}$}} q2;
output {\textcolor{highlight}{$r_{13}$}} q8;
output {\textcolor{highlight}{$r_{14}$}} q1;
\end{yquant}
\node[draw=highlight, thick, fill=highlight, fill opacity=.12, rounded corners=2pt, inner sep=2pt, fit=(m8-0)] (m8box) {};
\node[font=\scriptsize, text=highlight, align=center] (m8lab) at ([yshift=4mm]m8box.north |- current bounding box.north) {central readout, $r_8$};
\draw[-{Stealth[length=1.4mm]}, highlight, line width=.5pt] (m8lab.south) -- (m8box.north);
\node[draw=highlight, thick, fill=highlight, fill opacity=.10, rounded corners=2pt, inner sep=2.6pt, fit=(dexit-0) (dexit-6), label={[font=\footnotesize, text=highlight, yshift=2.5mm]above:$D$}] {};
\node[draw=highlight, dashed, fill=highlight, fill opacity=.06, rounded corners=2pt, inner sep=2.6pt, fit=(dentry-0) (dentry-6), label={[font=\footnotesize, text=highlight, yshift=2.5mm]above:$D^\dagger$}] {};
\end{tikzpicture}

\end{adjustbox}
\caption{The standalone $d=3$ double checking circuit. The entry layer $D^\dagger$ and the exit layer $D$ of $T$ gates are highlighted. Readouts are labelled by binary outcome $r_i$. 
}
\label{fig:d3circuit}
\end{figure}
\begin{figure}[p]\centering
\begin{adjustbox}{max width=.98\linewidth,max totalheight=.88\textheight}
\begin{tikzpicture}[font=\scriptsize,line width=.45pt]
\begin{yquant}[register/minimum height=1.1mm,register/minimum depth=1.1mm,register/separation=.6mm,operator/minimum width=4mm,operator/separation=2.2mm]
qubit {$24\;\ket+$} q24;
qubit {$25$} q25;
qubit {$31\;\ket+$} q31;
qubit {$32$} q32;
qubit {$37\;\ket+$} q37;
qubit {$36$} q36;
qubit {$41\;\ket+$} q41;
qubit {$40$} q40;
qubit {$26\;\ket+$} q26;
qubit {$19$} q19;
qubit {$33\;\ket+$} q33;
qubit {$34$} q34;
qubit {$39\;\ket+$} q39;
qubit {$38$} q38;
qubit {$20\;\ket+$} q20;
qubit {$27$} q27;
qubit {$28\;\ket+$} q28;
qubit {$21$} q21;
qubit {$22\;\ket+$} q22;
qubit {$29$} q29;
qubit {$30\;\ket+$} q30;
qubit {$23$} q23;
qubit {$18\;\ket+$} q18;
qubit {$17$} q17;
qubit {$12\;\ket+$} q12;
qubit {$11$} q11;
qubit {$8\;\ket+$} q8;
qubit {$13$} q13;
qubit {$14\;\ket+$} q14;
qubit {$9$} q9;
qubit {$10\;\ket+$} q10;
qubit {$15$} q15;
qubit {$6\;\ket+$} q6;
qubit {$7$} q7;
qubit {$2\;\ket+$} q2;
qubit {$3$} q3;
qubit {$1\;\ket+$} q1;
qubit {$0$} q0;
box {$T$} q36,q40,q19,q17,q11,q0;
box {$T^\dagger$} q25,q32,q34,q38,q27,q21,q29,q23,q13,q9,q15,q7,q3;
align q24,q25,q31,q32,q37,q36,q41,q40,q26,q19,q33,q34,q39,q38,q20,q27,q28,q21,q22,q29,q30,q23,q18,q17,q12,q11,q8,q13,q14,q9,q10,q15,q6,q7,q2,q3,q1,q0;
cnot q0 | q1;
cnot q3 | q2;
cnot q7 | q6;
cnot q13 | q8;
cnot q15 | q10;
cnot q11 | q12;
cnot q9 | q14;
cnot q17 | q18;
cnot q27 | q20;
cnot q29 | q22;
cnot q25 | q24;
cnot q19 | q26;
cnot q21 | q28;
cnot q23 | q30;
cnot q32 | q31;
cnot q34 | q33;
cnot q36 | q37;
cnot q38 | q39;
cnot q40 | q41;
align q24,q25,q31,q32,q37,q36,q41,q40,q26,q19,q33,q34,q39,q38,q20,q27,q28,q21,q22,q29,q30,q23,q18,q17,q12,q11,q8,q13,q14,q9,q10,q15,q6,q7,q2,q3,q1,q0;
cnot q1 | q3;
cnot q15 | q14;
cnot q30 | q29;
cnot q38 | q33;
cnot q41 | q37;
align q24,q25,q31,q32,q37,q36,q41,q40,q26,q19,q33,q34,q39,q38,q20,q27,q28,q21,q22,q29,q30,q23,q18,q17,q12,q11,q8,q13,q14,q9,q10,q15,q6,q7,q2,q3,q1,q0;
cnot q3 | q7;
cnot q14 | q13;
cnot q33 | q26;
cnot q29 | q28;
cnot q37 | q32;
align q24,q25,q31,q32,q37,q36,q41,q40,q26,q19,q33,q34,q39,q38,q20,q27,q28,q21,q22,q29,q30,q23,q18,q17,q12,q11,q8,q13,q14,q9,q10,q15,q6,q7,q2,q3,q1,q0;
cnot q13 | q12;
cnot q32 | q25;
cnot q28 | q27;
align q24,q25,q31,q32,q37,q36,q41,q40,q26,q19,q33,q34,q39,q38,q20,q27,q28,q21,q22,q29,q30,q23,q18,q17,q12,q11,q8,q13,q14,q9,q10,q15,q6,q7,q2,q3,q1,q0;
cnot q7 | q12;
cnot q27 | q26;
align q24,q25,q31,q32,q37,q36,q41,q40,q26,q19,q33,q34,q39,q38,q20,q27,q28,q21,q22,q29,q30,q23,q18,q17,q12,q11,q8,q13,q14,q9,q10,q15,q6,q7,q2,q3,q1,q0;
cnot q12 | q18;
cnot q26 | q25;
align q24,q25,q31,q32,q37,q36,q41,q40,q26,q19,q33,q34,q39,q38,q20,q27,q28,q21,q22,q29,q30,q23,q18,q17,q12,q11,q8,q13,q14,q9,q10,q15,q6,q7,q2,q3,q1,q0;
cnot q18 | q25;
align q24,q25,q31,q32,q37,q36,q41,q40,q26,q19,q33,q34,q39,q38,q20,q27,q28,q21,q22,q29,q30,q23,q18,q17,q12,q11,q8,q13,q14,q9,q10,q15,q6,q7,q2,q3,q1,q0;
[name=mc, type=qubit] measure {$X$} q25;
box {$R_X$} q25;
align q24,q25,q31,q32,q37,q36,q41,q40,q26,q19,q33,q34,q39,q38,q20,q27,q28,q21,q22,q29,q30,q23,q18,q17,q12,q11,q8,q13,q14,q9,q10,q15,q6,q7,q2,q3,q1,q0;
cnot q18 | q25;
align q24,q25,q31,q32,q37,q36,q41,q40,q26,q19,q33,q34,q39,q38,q20,q27,q28,q21,q22,q29,q30,q23,q18,q17,q12,q11,q8,q13,q14,q9,q10,q15,q6,q7,q2,q3,q1,q0;
cnot q12 | q18;
cnot q26 | q25;
align q24,q25,q31,q32,q37,q36,q41,q40,q26,q19,q33,q34,q39,q38,q20,q27,q28,q21,q22,q29,q30,q23,q18,q17,q12,q11,q8,q13,q14,q9,q10,q15,q6,q7,q2,q3,q1,q0;
cnot q7 | q12;
cnot q27 | q26;
align q24,q25,q31,q32,q37,q36,q41,q40,q26,q19,q33,q34,q39,q38,q20,q27,q28,q21,q22,q29,q30,q23,q18,q17,q12,q11,q8,q13,q14,q9,q10,q15,q6,q7,q2,q3,q1,q0;
cnot q13 | q12;
cnot q32 | q25;
cnot q28 | q27;
align q24,q25,q31,q32,q37,q36,q41,q40,q26,q19,q33,q34,q39,q38,q20,q27,q28,q21,q22,q29,q30,q23,q18,q17,q12,q11,q8,q13,q14,q9,q10,q15,q6,q7,q2,q3,q1,q0;
cnot q3 | q7;
cnot q14 | q13;
cnot q33 | q26;
cnot q29 | q28;
cnot q37 | q32;
align q24,q25,q31,q32,q37,q36,q41,q40,q26,q19,q33,q34,q39,q38,q20,q27,q28,q21,q22,q29,q30,q23,q18,q17,q12,q11,q8,q13,q14,q9,q10,q15,q6,q7,q2,q3,q1,q0;
cnot q1 | q3;
cnot q15 | q14;
cnot q30 | q29;
cnot q38 | q33;
cnot q41 | q37;
align q24,q25,q31,q32,q37,q36,q41,q40,q26,q19,q33,q34,q39,q38,q20,q27,q28,q21,q22,q29,q30,q23,q18,q17,q12,q11,q8,q13,q14,q9,q10,q15,q6,q7,q2,q3,q1,q0;
cnot q0 | q1;
cnot q3 | q2;
cnot q7 | q6;
cnot q13 | q8;
cnot q15 | q10;
cnot q11 | q12;
cnot q9 | q14;
cnot q17 | q18;
cnot q27 | q20;
cnot q29 | q22;
cnot q25 | q24;
cnot q19 | q26;
cnot q21 | q28;
cnot q23 | q30;
cnot q32 | q31;
cnot q34 | q33;
cnot q36 | q37;
cnot q38 | q39;
cnot q40 | q41;
align q24,q25,q31,q32,q37,q36,q41,q40,q26,q19,q33,q34,q39,q38,q20,q27,q28,q21,q22,q29,q30,q23,q18,q17,q12,q11,q8,q13,q14,q9,q10,q15,q6,q7,q2,q3,q1,q0;
box {$T$} q25,q32,q34,q38,q27,q21,q29,q23,q13,q9,q15,q7,q3;
box {$T^\dagger$} q36,q40,q19,q17,q11,q0;
align q24,q25,q31,q32,q37,q36,q41,q40,q26,q19,q33,q34,q39,q38,q20,q27,q28,q21,q22,q29,q30,q23,q18,q17,q12,q11,q8,q13,q14,q9,q10,q15,q6,q7,q2,q3,q1,q0;
measure {$X$} q20,q22,q8,q10,q2,q6,q24,q31,q39,q1,q12,q18,q37,q33,q14,q26,q28,q30,q41;
\end{yquant}
\node[draw=highlight, thick, fill=highlight, fill opacity=.12, rounded corners=2pt, inner sep=2pt, fit=(mc-0), label={[font=\scriptsize, text=highlight, yshift=3.6mm]above:central readout}] {};
\end{tikzpicture}

\end{adjustbox}
\caption{The $d=5$ double checking circuit, with its corrected $T$ pattern. Wires are ordered by the fold-layer components, and four sites idle in this window are not drawn. The central $X$ measurement on site 25 is highlighted.
}
\label{fig:d5circuit}
\end{figure}
